% ============================================================
%  CHAPTER 5 — Implementation
% ============================================================

\chapter{Implementation}
\label{ch:implementation}

\section{System Overview}

The firmware maps the software architecture described in
Chapter~\ref{ch:architecture} directly onto the ESP32's dual-core execution model.
At boot, \texttt{setup()} initialises LittleFS, loads NVS configuration, initialises the sensors, starts
both Wi-Fi interfaces, registers HTTP routes, and launches the Firebase FreeRTOS task .
After \texttt{setup()} returns, the Arduino \texttt{loop()} continues to execute, calling \texttt{updateSensors()} on every iteration and scheduling periodic
Wi-Fi background scans.

The \textbf{Backend Subsystem} is implemented as a Node.js service that acts as a secure proxy between the clinical dashboard and the Firebase infrastructure. It consumes data from Firestore for administrative operations and stores detection/alert logs. The server is initialized via a tRPC router which binds the service logic to the network layer, ensuring that every API call is strictly validated against the project's TypeScript interfaces. Authentication uses Firebase ID tokens at login and persists sessions via signed cookies.

\begin{figure}[H]
    \centering
    \includegraphics[width=\linewidth]{Backend Service Overview.png}
    \caption{Backend service and data integration overview showing tRPC/Firebase interaction.}
    \label{fig:backend_overview}
\end{figure}

The \textbf{Web Dashboard} is a high-fidelity React application that serves as the central command center for hospital staff. It implements a reactive state model where UI components automatically re-render in response to Firebase Realtime Database events. This tight coupling between the cloud database and the frontend components allows for a "zero-refresh" monitoring experience, which is critical for observing acute physiological changes. Threshold-based clinical alerts are computed on the client and trigger emails through a backend SMTP procedure.

% ============================================================
\section{Hardware Implementation}
% ============================================================

The hardware assembly consists of three components mounted on a common 3.3\,V power
rail:

\begin{itemize}
	\item \textbf{ESP32 development board}.
	\item \textbf{DHT22} connected to GPIO~4.
	\item \textbf{MAX30102 breakout board} connected to the ESP32's default
	      I\textsuperscript{2}C bus (SDA: GPIO~21, SCL: GPIO~22), powered at 3.3\,V.
\end{itemize}

All components are powered from the ESP32 development board's 3.3\,V regulated output.
The board is powered from USB 5\,V connected to a mains adapter in the hospital room.

% ============================================================
\section{Hardware Implementation Challenges}
% ============================================================

\subsection*{Camera Module: ESP-CAM Evaluation and Replacement}

During the initial design phase, an ESP32-CAM module was evaluated as an integrated solution for room monitoring and personnel identification via face recognition. The module was wired and tested within the existing hardware setup; however, it was ultimately rejected for this use case. The onboard OV2640 camera sensor produced images at an effective resolution and quality insufficient for reliable face detection: under standard indoor lighting conditions, captured frames exhibited low contrast and significant noise artifacts, making it impossible to extract discriminative facial descriptors for identification. As a result, the face recognition feature was migrated to an external USB camera connected directly to the monitoring terminal running the web dashboard, where image quality and frame rate are adequate for the face-api.js inference pipeline.

\subsection*{MAX30102 Pulse Oximeter: Pull-up Resistor Issue}

A hardware-level issue was encountered with the MAX30102 breakout board. In its factory configuration, the breakout board includes on-board pull-up resistors on the I\textsuperscript{2}C lines (SDA and SCL) and the interrupt line (INT). These resistors were found to be incompatible with the ESP32's I\textsuperscript{2}C bus voltage levels under the target operating conditions, causing communication failures and preventing reliable sensor readings. The resolution required physically removing the internal pull-up resistors from the breakout board and replacing them with external resistors of appropriate values connected directly between the 3.3\,V rail and the INT, SCL, and SDA pins. After this modification, I\textsuperscript{2}C communication was restored and the sensor operated correctly at 400\,kHz fast mode.

% ============================================================
\section{Hardware Schematics}
% ============================================================

Figure~\ref{fig:schematic} provides the complete wiring diagram of the assembled circuit.


\begin{figure}[H]
	\centering
	\includegraphics[width=\linewidth]{hardware_diagram.png}
	\caption{Hardware wiring diagram}
	\label{fig:schematic}
\end{figure}

% ============================================================
\section{Software Development Environment}
% ============================================================

Table~\ref{tab:dev_env} lists the complete development environment used for the
ESP32 embedded subsystem.

\begin{table}[H]
	\centering
	\caption{ESP32 software development environment.}
	\label{tab:dev_env}
	\begin{tabularx}{\linewidth}{lX}
		\toprule
		\textbf{Component} & \textbf{Details}                               \\
		\midrule
		Operating System   & Linux (Debian 13 trixie)                       \\
		ESP32 Arduino core & v3.3.7                                          \\
		Board target       & \texttt{esp32:esp32:esp32da} (ESP32-WROOM-32)  \\
		Flash tool         & \texttt{esptool} (Python); \texttt{mklittlefs} \\
		Serial monitor     & Custom \texttt{monitor.sh}                     \\
		\bottomrule
	\end{tabularx}
\end{table}

\begin{table}[H]
	\centering
	\caption{Full-stack software development environment.}
	\label{tab:backend_dev_env}
	\begin{tabularx}{\linewidth}{lX}
		\toprule
		\textbf{Component} & \textbf{Details}                               \\
		\midrule
		Language / Runtime & TypeScript / Node.js v20+                       \\
		Backend Framework  & tRPC v11 (Functional Router Architecture)      \\
		Frontend Framework & React v18 + Vite (SPA)                         \\
		Styling / UI       & TailwindCSS + Shadcn UI + Lucide Icons         \\
		Database           & Firebase Firestore, Realtime Database         \\
		Dependency Manager & PNPM / NPM                                     \\
		IDE                & VS Code with Prettier, ESLint                 \\
		\bottomrule
	\end{tabularx}
\end{table}

% ============================================================
\section{Software Implementation}
% ============================================================

\subsection{Embedded Subsystem}

\subsubsection{Firebase Initialisation on Core~0}

The \texttt{initFirebase()} function is called from inside
\texttt{firebaseUploadTask()}, which is pinned to Core~0 via
\texttt{xTaskCreatePinnedToCore(..., 0)}.

\subsubsection{Configuration Portal}

The configuration portal is served entirely from the LittleFS partition.
The \texttt{isAPConnected()} guard checks whether the incoming request arrived via
the AP subnet (\texttt{WiFi.softAPIP()}) and redirects any STA-side request to the
root path, preventing access from the hospital network.

\subsubsection{Wi-Fi Recovery}

All reconnection logic is consolidated in the \texttt{wifiConnect()} helper, which
always calls \texttt{WiFi.disconnect()} followed by
\texttt{WiFi.begin(ssid, password)}.

\subsection{Backend Implementation}
The backend implementation is modularized into \textbf{functional routers} (\texttt{people}, \texttt{events}, \texttt{alerts}, \texttt{ai}). The \texttt{patients} router is an alias of \texttt{people}, enabling a shared CRUD layer. Each router defines its procedures using Zod for runtime schema validation. Email delivery is implemented in a dedicated SMTP module and is invoked explicitly by the dashboard for threshold alerts. Authentication is implemented as a middleware layer that validates a signed session cookie derived from Firebase ID tokens.

Photo enrollment uses a storage adapter that falls back to base64 data URLs if
Firebase Storage is unavailable. This keeps the feature functional without a
paid storage plan but increases payload sizes in the database and client.

\subsection{Web Dashboard Implementation}
The dashboard's core is the \textbf{Real-time Monitoring Component}, which uses a custom hook to manage the Firebase RTDB listener lifecycle. Vitals are visualized using \textbf{AreaCharts} from Recharts, optimized for high-frequency updates without UI stutter. A secondary \textbf{AI Insights Panel} renders synthesized clinical conclusions by subscribing to the tRPC AI results. The personnel tracking system is implemented via a browser-based face-api.js pipeline, executing detection every 600\,ms and applying a 5-frame stability threshold before logging alerts.

% ============================================================
\section{Software Validation}
% ============================================================

\subsection{Embedded Subsystem Validation}

Validation was performed across three dimensions:

\textbf{Sensor output validation.}
DHT22 readings were cross-referenced against a calibrated reference thermometer.

\textbf{Firebase upload validation.}
The Firebase Realtime Database console was monitored in real time to confirm record
arrival, correct path structure (\texttt{/rooms/\{roomId\}} and
\texttt{/patients/\{patientId\}}), correct field names and data types, and
correct server-generated timestamps.

\textbf{Authentication validation.}
The serial monitor (debug build) was used to trace the complete authentication event
sequence: \textit{authenticating} $\to$ \textit{auth request sent} $\to$
\textit{auth response received} $\to$ \textit{token obtained}.
Error conditions (invalid API key, incorrect password) were deliberately induced
and confirmed to produce the expected HTTP~400 error codes via the callback.

\subsection{Backend Validation}
The backend includes a suite of \textbf{integration tests} using a tRPC caller. Test cases include:
\begin{itemize}
    \item \textbf{Auth Verification:} Confirming that requests without valid tokens return HTTP 401.
    \item \textbf{Schema Enforcement:} Ensuring that malformed JSON payloads are rejected by Zod validators.
	\item \textbf{Alert Dispatch:} Verifying that the SMTP relay procedure can be invoked and returns a success/failure response.
\end{itemize}

\subsection{Web Dashboard Validation}
UI validation focuses on \textbf{performance under load} and \textbf{clinical accuracy}.
\begin{itemize}
	\item \textbf{Stress Testing:} Assess UI responsiveness while rendering multiple chart streams and a concurrent camera loop.
	\item \textbf{Personnel Detection:} Validate that face recognition matches enrolled identities using the 0.6 confidence threshold and a 5-frame stability gate.
	\item \textbf{Responsive Design:} Validate layout integrity across Desktop, Tablet, and Mobile viewport sizes using browser tooling.
\end{itemize}

% ============================================================
\section{Embedded Implementation (ESP32)}
% ============================================================

\subsection{Memory Optimisation}

The LittleFS image partition is mounted with \texttt{LittleFS.begin(true)}, enabling automatic
formatting on first boot.
NVS storage uses the \texttt{Preferences} library with the \texttt{"settings"}
namespace, consuming at most 10 string entries at typical credential lengths.

\subsection{Latency Optimisation}

Delta filtering (Equations~\ref{eq:room_thresh} and~\ref{eq:patient_thresh}) reduces
the frequency of HTTPS requests to only those carrying new information, keeping the
per-upload round-trip time as the dominant latency term rather than queuing delays.

\section{AI Integration}
The AI prediction layer is integrated as a \textbf{server-side service} that processes data on-demand from the dashboard. This deployment strategy was chosen to leverage the Node.js server's computational resources, keeping the ESP32 firmware lightweight.
\begin{itemize}
	\item \textbf{Integration:} The forecasting logic is wrapped in a tRPC query that accepts a telemetry window (target length 20, minimum 5).
	\item \textbf{Optimization:} Z-score calculations and linear regressions are performed on in-memory arrays to minimize inference latency.
	\item \textbf{Validation:} Synthetic datasets are used during development; formal clinical validation remains an open requirement.
\end{itemize}
